\chapter{Modules, tensor products, and commutative algebra}\label{chap:M07.22}

All rings in this chapter are commutative and unital, all modules are
unital, and homomorphisms of rings preserve one. The zero ring is allowed;
its only unital module is zero. A subring in an integrality assertion has
the same identity as the ambient ring. These conventions are essential
for scalar extension and localization.

% Free modules and tensor products
\KEntry{K07.22.free_direct_sums}{Free modules and arbitrary direct sums}
For a set $E$, let $R^{(E)}$ be the module of functions $E\to R$ with
finite support, with pointwise operations. Write $\delta_e$ for the
function supported at $e$ with value one. Every element is the finite sum
$\sum_e a_e\delta_e$, uniquely. Consequently each function $h:E\to M$
extends uniquely to an $R$-linear map
$R^{(E)}\to M$, namely $\sum a_e\delta_e\mapsto\sum a_eh(e)$.
This proves both existence and the defining universal property of a free
module, without assuming that $E$ is finite.

For a family $(M_i)_{i\in E}$ define $\bigoplus_iM_i$ to be the
finite-support tuples in $\prod_iM_i$. The coordinate inclusions
$\iota_i:M_i\to\bigoplus_iM_i$ have the following property: $R$-linear maps
$h_i:M_i\to N$ extend uniquely to
$h((m_i))=\sum_i h_i(m_i)$. The sum is finite. This follows by writing
$(m_i)=\sum_i\iota_i(m_i)$; the same equality proves uniqueness and
linearity. There is no finite-support restriction in the product
$\prod_iM_i$, so this universal property does not describe the product
for an infinite index set. Empty direct sums and free modules on the
empty set are zero.

\KEntry{K07.22.tensor_construction}{The balanced tensor quotient}
A map $b:M\times N\to A$ to an abelian group is balanced if it is additive
in each argument and $b(rm,n)=b(m,rn)$ for all $r\in R$. Start with the
free abelian group on the set $M\times N$, whose existence is the preceding
construction with $R=\Z$. Divide by the subgroup generated by
\[
 [m+m',n]-[m,n]-[m',n],\quad
 [m,n+n']-[m,n]-[m,n'],\quad [rm,n]-[m,rn].
\]
Call the quotient $M\otimes_RN$ and the class of $[m,n]$ the elementary
tensor $m\otimes n$. A homomorphism from the free group factors through
this quotient exactly when it kills the displayed generators.
It follows that every balanced map factors uniquely through an additive
map $M\otimes_RN\to A$. In particular elementary tensors generate the
underlying abelian group.

This group is an $R$-module by $r(m\otimes n)=(rm)\otimes n$.
To justify well-definedness, the map $[m,n]\mapsto[rm,n]$ takes the two
additivity relations into relations. It takes a balancing relation for
$s$ into
$[rsm,n]-[rm,sn]$, which is a balancing relation since $rs=sr$.
Thus it descends. Additivity in $r$, multiplication of scalars and the
identity axiom hold on elementary tensors by the defining relations;
they hold on the whole group because those tensors generate it.
For an $R$-module $P$, an $R$-bilinear map $M\times N\to P$ is balanced,
and its induced additive map is $R$-linear on elementary tensors and hence
everywhere. We have proved the universal property
\[
 \Hom_R(M\otimes_RN,P)\ \simeq\
 \{\text{$R$-bilinear maps }M\times N\to P\}.
\]
Merely being balanced into an abelian group is a different formulation:
there is then no scalar-linearity assertion about the target.
This construction extends the field-only tensor interface
\texttt{K05.15.tensor\_existence} and
\texttt{K05.15.tensor\_universal}; it does not assume a basis for $M$ or $N$.

\KEntry{K07.22.tensor_functor}{Maps and canonical tensor isomorphisms}
For linear maps $f:M\to M'$ and $g:N\to N'$, bilinearity constructs the
unique map $f\otimes g$ taking $m\otimes n$ to $f(m)\otimes g(n)$.
The equations
\[
 (f'f)\otimes(g'g)=(f'\otimes g')(f\otimes g),\qquad
 \id_M\otimes\id_N=\id_{M\otimes N}
\]
hold because the two sides agree on every elementary tensor.
This argument also proves additivity in each map.
It is a useful general principle here: two linear maps out of a tensor
product are equal once their values on elementary tensors are equal.

There are natural isomorphisms
\[
 M\otimes N\simeq N\otimes M,\quad
 (M\otimes N)\otimes P\simeq M\otimes(N\otimes P),\quad
 R\otimes M\simeq M.
\]
For the first use $m\otimes n\mapsto n\otimes m$, whose square is the
identity. For the second, for fixed $p$ the bilinear map
$(m,n)\mapsto m\otimes(n\otimes p)$ induces a linear map on $M\otimes N$.
Its dependence on $p$ is linear: the equations can be checked first on
$m\otimes n$ and then on their sums. The resulting bilinear map
$(M\otimes N)\times P\to M\otimes(N\otimes P)$ induces the asserted map.
Interchanging the parentheses constructs its inverse, and the
composites fix $(m\otimes n)\otimes p$ and $m\otimes(n\otimes p)$.
These elements generate the respective iterated tensor products.
For the last map use $r\otimes m\mapsto rm$, with inverse $m\mapsto1\otimes m$;
the balancing relation gives $1\otimes rm=r\otimes m$.
This proves the unit assertion even for the zero ring.

\KEntry{K07.22.tensor_direct_sums}{Tensor products distribute over direct sums}
For every family $(M_i)$ there is a canonical isomorphism
\[
 \left(\bigoplus_iM_i\right)\otimes N
 \ \simeq\ \bigoplus_i(M_i\otimes N).
\]
The map on a pair $((m_i),n)$ is the tuple $(m_i\otimes n)_i$.
It has finite support, is bilinear, and therefore induces the forward
map. In the reverse direction the maps
$\iota_i\otimes\id_N$ extend by the direct-sum universal property.
The reverse composite takes $(m_i)\otimes n$ to
$\sum_i\iota_i(m_i)\otimes n=(m_i)\otimes n$.
The other composite fixes each included $m_i\otimes n$.
Generation proves both identities. This proof works for infinite
families because every sum used has finite support. In particular
$R^{(E)}\otimes N\simeq N^{(E)}$ and
$R^a\otimes R^b\simeq R^{ab}$ for finite nonnegative $a,b$.
For nonzero commutative $R$ this is multiplication of finite free ranks:
modulo a maximal ideal, any free-basis isomorphism becomes an isomorphism
of vector spaces, whose dimensions are equal. Over the zero ring
every such module is zero and its displayed basis size is not an
intrinsic rank. No corresponding assertion about arbitrary infinite
products follows from this argument.

\KEntry{K07.22.tensor_quotients}{Quotient tensors and cyclic examples}
\begin{proposition}\label{prop:tensor-quotient}
For submodules $U\subseteq M$ and $V\subseteq N$, let $J\subseteq M\otimes N$
be the submodule generated by $u\otimes n$ and $m\otimes v$, with
$u\in U,v\in V,m\in M,n\in N$. Then
\[
 (M/U)\otimes_R(N/V)\ \simeq\ (M\otimes_RN)/J.
\]
\end{proposition}
\begin{proof}
The bilinear map $(m,n)\mapsto(m+U)\otimes(n+V)$ kills the generators of
$J$, so induces a map from the right-hand quotient. Conversely set
$(m+U,n+V)\mapsto m\otimes n+J$. Replacing $m$ by $m+u$ changes the
value by $u\otimes n\in J$, and replacing $n$ by $n+v$ changes it by
$(m+u)\otimes v\in J$. The map is well-defined and bilinear.
It induces the reverse map. Their composites fix all the displayed
generators, which proves they are inverses.
\end{proof}
The definition of $J$ uses images of tensor maps. It does not assert that
$U\otimes N\to M\otimes N$ is injective.
Taking $M=R,U=I,V=0$ and using the tensor unit gives
\begin{equation}\label{eq:quotient-tensor}
 (R/I)\otimes_RN\simeq N/IN,
\end{equation}
where $IN$ is the submodule of finite sums of $in$ with $i\in I,n\in N$.
Indeed $r\otimes n\mapsto rn$ takes $I$-tensors exactly onto $IN$.
For integers $m,n$, this gives
\[
 \Z/(m)\otimes_\Z\Z/(n)\simeq\Z/(m,n)
       \simeq\Z/(\gcd(m,n)).
\]
The last equality uses the integer Bezout identity, with the nonnegative
gcd and $\gcd(0,0)=0$. Thus a zero modulus means the infinite cyclic
module; a unit modulus gives zero.

\KEntry{K07.22.tensor_right_exact}{Right exactness and failure of left exactness}
\begin{theorem}\label{thm:tensor-right-exact}
If $A\xrightarrow fB\xrightarrow gC\to0$ is exact, then
\[
 A\otimes N\xrightarrow{f\otimes\id}B\otimes N
 \xrightarrow{g\otimes\id}C\otimes N\to0
\]
is exact.
\end{theorem}
\begin{proof}
Exactness means that $g$ is onto and $\ker g=\im f$. The map
$B/\im f\to C$ induced by $g$ is a bijective linear map: surjectivity
comes from $g$, and a zero image means the representative lies in its
kernel. Apply Proposition~\ref{prop:tensor-quotient} with $U=\im f,V=0$.
Its submodule $J$ is generated by $f(a)\otimes n$ and equals the image
of $f\otimes\id$: every generator is the image of $a\otimes n$, and
every tensor in the source is a sum of such elements. The quotient map
is onto and has kernel $J$, proving both required exactness assertions.
\end{proof}
Left exactness can fail. The injective map $\Z\to\Z$, $x\mapsto2x$,
becomes the zero map on $\Z\otimes_\Z\Z/(2)\simeq\Z/(2)$ after tensoring.
Its source is nonzero, so it is not injective. This is also a concrete
reason not to regard tensor products of submodules as submodules
without an additional hypothesis such as flatness.

\KEntry{K07.22.scalar_extension}{Scalar extension and its universal property}
Let $\varphi:R\to S$ be a unital homomorphism of commutative rings.
Regard $S$ as an $R$-module and make $S\otimes_RM$ an $S$-module by
$s'(s\otimes m)=s's\otimes m$. The defining tensor relations are
preserved, so this action is well-defined. For an $S$-module $N$,
\[
 \Hom_S(S\otimes_RM,N)\simeq\Hom_R(M,N),
\]
where on the right scalars act through $\varphi$.
The forward map sends $F$ to $m\mapsto F(1\otimes m)$.
For an $R$-linear $h$, the map $(s,m)\mapsto s\,h(m)$ is balanced and
biadditive, and hence induces a map $s\otimes m\mapsto s\,h(m)$.
It is $S$-linear on elementary tensors, so everywhere.
The two constructions are inverse because $S$-linearity determines the
value at every $s\otimes m$ from the value at $1\otimes m$.

The map $M\to S\otimes_RM$, $m\mapsto1\otimes m$, need not be injective:
if $M$ is a finite torsion abelian group, $\Q\otimes_\Z M=0$.
Indeed an elementary tensor $q\otimes m$, with $dm=0,d\ne0$, equals
$(q/d)\otimes dm=0$. If $F\subseteq L$ are fields and $(e_i)$ is an
$F$-basis of $M$, the direct-sum formula gives
$L\otimes_FM\simeq\bigoplus_iL$ with basis $1\otimes e_i$.
In these coordinates the natural map takes the $F$-coordinates to the
same coordinates in $L$, so is injective and preserves finite dimension
after extending scalars. The field case recovers
\texttt{K05.15.tensor\_maps}, while the torsion example has no field analogue.

\KEntry{K07.22.complex_tensor_example}{A scalar-extension algebra that splits}
For commutative $R$-algebras $A,B$, the tensor module has an algebra
structure
\[
 (a\otimes b)(a'\otimes b')=aa'\otimes bb',\qquad1=1\otimes1.
\]
Here is a construction, rather than an assumption about representatives.
For fixed $a',b'$, the displayed expression is bilinear in $a,b$;
the resulting linear map in $A\otimes B$ is bilinear in $a',b'$,
as can be checked on its elementary generators. Applying the universal
property again yields a bilinear multiplication on two copies of the
tensor module. Associativity, commutativity, the unit property and
distributivity are equations of multilinear maps; on elementary
generators they follow from those laws in $A,B$, and expanding finite
sums proves them everywhere. The maps from $A$ and $B$ are unital
algebra maps with commuting images.

As a $\C$-algebra through the first factor there is an isomorphism
\[
 \C\otimes_\R\C\longrightarrow\C\times\C,\qquad
 z\otimes w\longmapsto(zw,z\overline w).
\]
The formula is $\R$-balanced, and its products and unit have the asserted
images, so it induces a $\C$-algebra homomorphism.
Writing $w=a+bi$ with $a,b\in\R$ shows that every tensor is uniquely of
the form $u\otimes1+v\otimes i$: existence follows by expansion, and
the proposed map sends this expression to $(u+iv,u-iv)$.
These two equations determine
$u=(x+y)/2,v=(x-y)/(2i)$ for any pair $(x,y)$.
They prove surjectivity, uniqueness of the expression, and injectivity.
The two nonzero orthogonal idempotents are
$\tfrac12(1\otimes1-i\otimes i)$ and
$\tfrac12(1\otimes1+i\otimes i)$, mapping to $(1,0)$ and $(0,1)$.
Thus a tensor product of fields need not be a domain. Only the real
basis $1,i$ and the relation $i^2=-1$ are used here.

% Finiteness conditions
\KEntry{K07.22.noetherian_equivalences}{Noetherian modules and the ascending-chain condition}
\begin{theorem}\label{thm:noetherian-equivalences}
For an $R$-module $M$ the following are equivalent:
every ascending chain of submodules stabilizes; every submodule of $M$
is finitely generated; every nonempty collection of submodules has a
maximal member under inclusion. Such a module is called Noetherian.
A ring is Noetherian when its regular module is Noetherian, or,
equivalently, every ideal is finitely generated.
\end{theorem}
\begin{proof}
Assume the ascending-chain condition. If a submodule $U$ were not
finitely generated, choose $u_1\in U$ and then
$u_{j+1}\in U\setminus\langle u_1,\ldots,u_j\rangle$.
This would give a strictly ascending chain, so no such $U$ exists.
Conversely the union of an ascending chain $U_1\subseteq U_2\subseteq\cdots$
is a submodule. Finitely many generators of the union all belong to
one $U_j$, so the chain stabilizes there. If a nonempty collection has
no maximal member, repeatedly choose a strictly larger member; the
resulting chain contradicts the condition. Finally if every such
collection has a maximal member, apply this to the members of an
ascending chain. A maximal $U_j$ must contain all later members.
These arguments use the usual choice principle for the successive
selections. Ideals are exactly submodules of the regular module.
\end{proof}

\KEntry{K07.22.noetherian_finite_modules}{Finite modules, submodules, and finite presentations}
\begin{proposition}\label{prop:noetherian-modules}
Submodules and quotients of a Noetherian module are Noetherian.
In an exact sequence $0\to U\to M\to V\to0$, if $U$ and $V$ are
Noetherian, then $M$ is Noetherian. Consequently every finitely generated
module over a Noetherian ring is Noetherian and finitely presented.
\end{proposition}
\begin{proof}
An ascending chain in a submodule is one in $M$. A chain in a quotient
lifts by inverse image to one in $M$, so it also stabilizes.
For the middle assertion, let $W\subseteq M$. Its intersection with
$U$ has finitely many generators $u_1,\ldots,u_s$. Its image in $V$
has finitely many generators, which lift to $w_1,\ldots,w_t\in W$.
Given $w\in W$, subtract a linear combination of the $w_j$ with the
same image. The remainder is in $W\cap U$ and hence is a combination
of the $u_i$. This proves that $W$ is finitely generated; use
Theorem~\ref{thm:noetherian-equivalences}.
Induction using $0\to R^{n-1}\to R^n\to R\to0$ makes every finite free
module Noetherian. A finitely generated $M$ is a quotient of some $R^n$,
so is Noetherian. The kernel of this quotient is a finitely generated
submodule of $R^n$; choosing generators gives an exact presentation
$R^s\to R^n\to M\to0$ with finite $s,n$.
\end{proof}
No freeness of the submodule is asserted. Finite generation over a
non-Noetherian ring alone does not yield these conclusions.

\KEntry{K07.22.hilbert_basis}{The Hilbert basis theorem}
\begin{theorem}\label{thm:hilbert-basis}
If $R$ is Noetherian, then $R[X_1,\ldots,X_n]$ is Noetherian for every
finite $n$. Every finitely generated commutative $R$-algebra is therefore
Noetherian.
\end{theorem}
\begin{proof}
It suffices to prove the one-variable assertion and iterate.
Let $I\subseteq R[X]$ be an ideal. The set consisting of zero and all
leading coefficients of its nonzero polynomials is an ideal $L$ of $R$.
For addition, multiply two witnessing polynomials by powers of $X$ to
align their degrees; if the sum of their leading coefficients is nonzero
it is the leading coefficient of their sum, and otherwise it is zero.
For multiplication by $r$, the leading coefficient of $rf$ is $ra$
when $ra\ne0$, and zero is already admitted when $ra=0$.
This argument allows zero divisors in $R$.
If $L=0$, there is no nonzero polynomial in $I$, so $I=0$.
Otherwise choose generators $a_1,\ldots,a_t$ of $L$, with witnessing
$f_i\in I$ of degrees $d_i$, and let $N=\max_i d_i$.

The polynomials of degree less than $N$, including zero, form a finite
free $R$-module $P_N$ with basis $1,X,\ldots,X^{N-1}$; when $N=0$ it
is zero. By Proposition~\ref{prop:noetherian-modules},
$I\cap P_N$ has module generators $h_1,\ldots,h_s$.
We claim that the $f_i,h_j$ generate $I$ as an ideal. For nonzero
$f\in I$ of degree $d\ge N$, write its leading coefficient as
$\sum_i r_ia_i$. Then
\[
 f-\sum_i r_iX^{d-d_i}f_i
\]
is zero or has smaller degree, and remains in $I$.
Induction on the nonnegative degree reduces to degree less than $N$,
where the remainder is an $R$-linear combination of the $h_j$.
Thus $I$ is finitely generated. Finite-variable induction proves the
first claim. A finite algebra generating list gives a surjection
$R[X_1,\ldots,X_n]\to A$; its quotient is a Noetherian module by the
quotient case of Proposition~\ref{prop:noetherian-modules}.
The submodules of this quotient module are precisely the ideals of
$A$, since scalar multiplication factors through the surjection and
every scalar of $A$ has a preimage. Thus $A$ is a Noetherian ring.
\end{proof}
The number of variables must be finite: in $k[X_1,X_2,\ldots]$ the ideals
$(X_1)\subsetneq(X_1,X_2)\subsetneq\cdots$ are strictly increasing,
since setting the earlier variables to zero leaves $X_{j+1}$ nonzero.
Finite algebra generation here does not mean finite module generation:
$k[X]$ is generated by $X$ as a $k$-algebra but has the infinite
$k$-linearly independent list $1,X,X^2,\ldots$.

% Localization
\KEntry{K07.22.localization_construction}{Localization with zero divisors}
Let $S\subseteq R$ contain $1$ and be closed under multiplication.
We temporarily allow $0\in S$. On $R\times S$ put
\[
 (a,s)\sim(b,t)\quad\Longleftrightarrow\quad
 \text{there is $u\in S$ with }u(ta-sb)=0.
\]
This extra multiplier is essential when there are zero divisors.
The relation is reflexive and symmetric. For transitivity, if
$u(ta-sb)=0$ and $v(wb-tc)=0$, the identity
$w(ta-sb)+s(wb-tc)=t(wa-sc)$ shows that $uvt(wa-sc)=0$;
the multiplier $uvt$ belongs to $S$.
Write $a/s$ for an equivalence class and define
\[
 \frac as+\frac bt=\frac{at+bs}{st},\qquad
 \frac as\frac bt=\frac{ab}{st},\qquad
 -\frac as=\frac{-a}s.
\]
For completeness, if $u(s'a-sa')=0$ and $v(t'b-tb')=0$, the difference
of cross products for the two sums is
\[
 tt'(s'a-sa')+ss'(t'b-tb'),
\]
and for the two products it is
\[
 bt'(s'a-sa')+a's(t'b-tb').
\]
Multiplication by $uv$ kills both. Negation is immediate.
Thus these operations are well-defined.
Associativity and commutativity of addition and multiplication, and
distributivity, follow by expanding numerators over common product
denominators: both sides have identical numerators in $R$.
The identities are $0/1$ and $1/1$, and
$(a/s)+(-a/s)=0/s^2=0/1$ because the cross difference is zero.
This constructs the ring $S^{-1}R$, including its operations and laws.
The map $\lambda:R\to S^{-1}R$, $a\mapsto a/1$, is a unital homomorphism.
Every $\lambda(s)$ is a unit with inverse $1/s$.

\KEntry{K07.22.localization_universal}{Universal property, kernel, and exceptional cases}
\begin{proposition}\label{prop:localization-universal}
Every unital homomorphism $f:R\to A$ for which $f(s)$ is a unit for
all $s\in S$ factors uniquely as a unital homomorphism
$\widetilde f:S^{-1}R\to A$, with
$\widetilde f(a/s)=f(a)f(s)^{-1}$. Moreover
\[
 \ker\lambda=\{a:\text{some }s\in S\text{ satisfies }sa=0\}.
\]
The localization is the zero ring exactly when $0\in S$.
The map $\lambda$ is injective exactly when multiplication by each
$s\in S$ is injective on $R$.
\end{proposition}
\begin{proof}
If $u(ta-sb)=0$, applying $f$ and cancelling the unit $f(u)$ gives
$f(t)f(a)=f(s)f(b)$. Multiplication by the inverse units proves that
the proposed images agree. The fraction formulas prove additivity,
multiplicativity and preservation of one. Every fraction equals
$\lambda(a)\lambda(s)^{-1}$, which proves uniqueness.
The equivalence $a/1=0/1$ is precisely the existence of $s$ with $sa=0$.
In particular $1/1=0/1$ exactly when an element of $S$ is zero.
For injectivity, the displayed kernel is zero exactly when
$sa=0$ for $s\in S$ implies $a=0$. This is the asserted injectivity
of every multiplication map.
\end{proof}
Usually one imposes $0\notin S$ to obtain a nonzero localization.
When the original ring is zero, every possible $S$ contains zero.
For a domain and $S=R\setminus\{0\}$, this recovers the fraction field
and its universal property from Chapter~\ref{chap:M07.08}.
For a ring with zero divisors localization may kill nonzero elements:
in $R=k[x,y]/(xy)$, localizing at powers of $x$ kills $y$.
Here $y\ne0$, as evaluation $x\mapsto0$ maps it to the nonzero
polynomial $y$ in $k[y]$, whereas $xy=0$.
For $S=\{1,f,f^2,\ldots\}$ there is a useful presentation
\[
 R_f\simeq R[T]/(fT-1).
\]
Indeed evaluation $T\mapsto1/f$ gives a homomorphism from the quotient
to $R_f$, and in the quotient the image of $f$ is invertible with inverse
the image of $T$. The universal property supplies the reverse map.
Both composites fix the images of $R$ and the indicated inverse
generator, so they are identities. This remains true for nilpotent $f$:
then $0\in S$ and both rings are zero.

\KEntry{K07.22.localization_ideals}{Localized ideals and prime correspondence}
\begin{theorem}\label{thm:localization-primes}
For an ideal $I\subseteq R$, its extension is
$S^{-1}I=\{a/s:a\in I,s\in S\}$ and its contraction is the saturation
\[
 \lambda^{-1}(S^{-1}I)=\{r:\text{some }s\in S\text{ has }sr\in I\}.
\]
Every ideal of $S^{-1}R$ is the extension of its contraction.
Extension and contraction restrict to inverse inclusion-preserving
bijections between prime ideals of $R$ disjoint from $S$ and prime
ideals of $S^{-1}R$.
\end{theorem}
\begin{proof}
The stated set of fractions is an ideal by the fraction formulas;
it contains $\lambda(I)$ and each of its fractions is
$\lambda(a)\lambda(s)^{-1}$, so it is exactly the generated ideal.
If $r/1=a/t$ with $a\in I$, there is $u\in S$ with $u(tr-a)=0$,
so $ut\,r=ua\in I$. Conversely $sr\in I$ gives
$r/1=(sr)/s$ in the extension. This proves saturation.
For any ideal $J$ of the localization, $a/s\in J$ implies
$a/1=(s/1)(a/s)\in J$; conversely $a/1\in J$ implies $a/s\in J$.
These two implications prove the extension-contraction assertion.
Also $S^{-1}I$ is the whole ring exactly when $I\cap S\ne\varnothing$.
One direction follows because an element of the intersection becomes
a unit. For the other, if $1=a/s$ with $a\in I$, a witness
$u(s-a)=0$ gives $us=ua\in I\cap S$.

Let $P$ be prime and disjoint from $S$. Its extension is proper by this
last criterion, and its contraction equals $P$: if $sr\in P$,
then $s\notin P$ forces $r\in P$.
If $(a/s)(b/t)\in S^{-1}P$, saturation applied to the numerator says
that some $u\in S$ has $uab\in P$; since $u\notin P$, primeness gives
$a\in P$ or $b\in P$. Hence $S^{-1}P$ is prime.
Conversely the inverse image of a prime ideal $Q$ of $S^{-1}R$ is prime
under a unital homomorphism: a product mapping into $Q$ has a factor
mapping into it, and one does not map into a proper ideal.
Its inverse image misses $S$, whose images are units.
The preceding extension-contraction equalities now prove that the two
operations are inverse. Inclusions are preserved from their definitions.
\end{proof}
If $0\in S$, both sets of prime ideals are empty, consistently with the
zero localization. No domain or Noetherian hypothesis entered the proof.

\KEntry{K07.22.local_rings}{Localization at a prime and local rings}
A nonzero commutative ring is local if it has a unique maximal ideal.
For a prime $P\subset R$, put $R_P=(R\setminus P)^{-1}R$.
This is nonzero, and its nonunits are exactly $PR_P$.
Indeed a fraction $a/s$ with $a\notin P$ has inverse $s/a$.
A fraction with numerator in $P$ lies in the proper ideal $PR_P$ and
cannot be a unit. Every proper ideal contains only nonunits, so every
proper ideal of $R_P$ is contained in $PR_P$.
This also proves that $PR_P$ itself is maximal and is the unique
maximal ideal. Thus $R_P/PR_P$ is a field: if the class of $x$ is nonzero,
the ideal it generates together with $PR_P$ is the whole ring, yielding
an inverse modulo $PR_P$. This argument is the usual maximal-ideal
quotient criterion, with its inverse construction made explicit.
The map $R\to R_P$ need not be injective.

% Integrality
\KEntry{K07.22.determinant_trick}{Integral elements and the determinant trick}
An element $x$ of a commutative $R$-algebra $B$ is integral over $R$ if
it satisfies a monic equation
$x^n+a_{n-1}x^{n-1}+\cdots+a_0=0$ with $n\ge1$ and $a_i\in R$.
When $R$ is a subring we identify coefficients with their images.
\begin{theorem}\label{thm:integral-finite}
For $x\in B$, the following are equivalent:
$x$ is integral over $R$; $R[x]$ is a finitely generated $R$-module;
$x$ belongs to a subalgebra of $B$ that is a finitely generated
$R$-module.
\end{theorem}
\begin{proof}
A monic equation expresses $x^n$ as a combination of
$1,x,\ldots,x^{n-1}$. Multiplying the equation by powers of $x$ and
inducting reduces every higher power, so these finitely many powers
span $R[x]$. This proves the first implication, and the second is
immediate by using $R[x]$.
For the last implication let $C\subseteq B$ be such a finite subalgebra
containing $x$, with module generators $v_1,\ldots,v_t$.
If $C$ is the zero ring then $B$ is zero because its identity belongs
to $C$, and the equation $x=0$ suffices.
Otherwise choose a nonempty generating list.
For each $i$ write $xv_i=\sum_j a_{ij}v_j$ with $a_{ij}\in R$.
For the column $v=(v_1,\ldots,v_t)^{\mathsf T}$ this is
$(xI-A)v=0$ over the commutative ring $B$.
Multiplying by the adjugate yields
$\det(xI-A)v_i=0$ for every $i$. The required commutative-ring
adjugate identity was proved by cofactor expansion in
Lemma~\ref{lem:commutative-ring-adjugate}; the field-only D05 statement
is not being used with broader hypotheses. Thus no division or
basis for $C$ is required. Since $1$ is a linear combination of the
$v_i$, it follows that $\det(xI-A)=0$ in $B$.
The polynomial $\det(TI-A)\in R[T]$ is monic of degree $t$:
the identity-permutation product supplies $T^t$, and every other
determinant term has smaller degree. This proves integrality.
\end{proof}
The same argument proves a useful determinant trick: if an element
$x$ of a commutative overring stabilizes a finitely generated submodule
$M$ of that ring, then a monic polynomial in $x$ annihilates $M$.
It gives a monic equation for $x$ itself when, for example, $1\in M$,
or the action on $M$ of the subring $R[x]$ of the overring is faithful.
Faithfulness merely as an $R$-module would not suffice.
Mere finite generation of an
unrelated module is insufficient.

\KEntry{K07.22.integral_closure_transitivity}{Finite integral algebras and transitivity}
\begin{theorem}\label{thm:integrality-transitive}
Elements integral over $R$ in a commutative $R$-algebra form a subalgebra.
An algebra generated by finitely many integral elements is a finite
$R$-module, and every finite $R$-algebra is integral.
If $R\subseteq B\subseteq C$, $B$ is integral over $R$, and $C$ is
integral over $B$, then $C$ is integral over $R$.
\end{theorem}
\begin{proof}
If $x_1,\ldots,x_t$ have monic equations over $R$ of degrees
$n_1,\ldots,n_t$, reduction of powers in each variable shows that
$R[x_1,\ldots,x_t]$ is spanned by
$x_1^{e_1}\cdots x_t^{e_t}$ with $0\le e_i<n_i$.
The reductions commute because the algebra is commutative.
Thus it is a finite module, and every one of its elements is integral
by Theorem~\ref{thm:integral-finite}. In particular sums, differences
and products of integral elements are integral, as are the elements
of $R$. A finite algebra is itself the finite subalgebra required
in that theorem, so all its elements are integral.

For transitivity fix $z\in C$ and choose its monic equation over $B$.
There are finitely many coefficients $b_1,\ldots,b_t$.
They are integral over $R$, so $A=R[b_1,\ldots,b_t]$ is a finite
$R$-module. The equation makes $A[z]$ a finite $A$-module.
If $a_i$ generate $A$ over $R$ and $w_j$ generate $A[z]$ over $A$,
then the finite list $a_iw_j$ generates $A[z]$ over $R$ by expanding
each coefficient. Thus $z$ belongs to a finite $R$-subalgebra of $C$
and is integral over $R$. This avoids any assertion that submodules
of finite modules are automatically finite over an arbitrary ring.
\end{proof}
It follows that a finitely generated algebra is finite as a module
exactly when all its algebra generators are integral.
An integral algebra without a finite algebra generating list need
not be a finite module. For example an algebraic closure of $\Q$
is integral over $\Q$ but has unbounded finite subextension degrees:
for every positive integer $n$, $X^n-2$ is Eisenstein at $2$, so its root has
degree $n$. A fixed finite vector-space dimension cannot accommodate
these subspaces.

\KEntry{K07.22.integrally_closed_ufd}{UFDs are integrally closed}
A domain $R$ is integrally closed if every element of $\Frac(R)$
integral over $R$ lies in $R$. Every UFD has this property.
Indeed write a fraction $r/s$ in reduced form with $s\ne0$, so no prime
divides both numerator and denominator. A monic equation of degree
$n$ gives, after multiplying by $s^n$,
\[
 r^n+a_{n-1}r^{n-1}s+\cdots+a_0s^n=0.
\]
Thus $s$ divides $r^n$. Any prime factor of a nonunit $s$ would divide
$r$ by primeness, contradicting reduction. Hence $s$ is a unit.
This includes $r=0$ by writing it as $0/1$.
In particular $1/2\in\Q$ is not integral over $\Z$, even though it is
algebraic over the fraction field $\Q$. Integrality over a ring uses a
monic equation with coefficients in that ring, a stronger condition
than an equation over its fraction field.

\KEntry{K07.22.lying_over}{The basic lying-over theorem}
\begin{theorem}\label{thm:lying-over}
For an integral extension of commutative rings $R\subseteq B$ and every
prime ideal $P$ of $R$, there is a prime ideal $Q$ of $B$ with
$Q\cap R=P$.
\end{theorem}
\begin{proof}
We first prove the field criterion used below. If domains $A\subseteq D$
form an integral extension and $D$ is a field, then for $0\ne a\in A$
the element $a^{-1}\in D$ is integral. Multiplying a monic equation
for $a^{-1}$ of degree $n$ by $a^{n-1}$ expresses $a^{-1}$ as an
element of $A$. Thus $A$ is a field. Conversely if $A$ is a field,
every nonzero $d\in D$ has a monic equation of least positive degree.
Its constant coefficient is nonzero, since otherwise cancellation
of $d$ in the domain would give a smaller monic equation.
Solving this equation for $d^{-1}$ puts that inverse in $A[d]$,
so $D$ is a field.
Consequently the contraction of a maximal ideal in any integral
extension is maximal: quotienting gives an integral extension of a
domain by a field, and the first direction of the criterion applies.
Integrality passes to these quotients simply by reducing each monic
equation.

Set $S=R\setminus P$. The map $S^{-1}R\to S^{-1}B$ induced by the
inclusion is injective. Indeed if $r/s$ maps to zero, some $u\in S$
satisfies $ur=0$ in $B$, and hence in the subring $R$, so $r/s=0$
already in $S^{-1}R$. The extension of localizations is integral:
if $b^n+a_{n-1}b^{n-1}+\cdots+a_0=0$, then $b/s$ satisfies the
monic equation with coefficients
$a_{n-1}/s,a_{n-2}/s^2,\ldots,a_0/s^n$.
The ring $R_P$ is nonzero; its injective image makes $S^{-1}B$ nonzero.

Choose a maximal ideal $M$ of $S^{-1}B$.
Such an ideal exists in every nonzero unital commutative ring:
partially order the proper ideals, note that the union of a chain is
an ideal and is still proper because it cannot contain one, and apply
Zorn's lemma. The contraction of $M$ to $R_P$ is maximal by the field
criterion, so it is the unique maximal ideal $PR_P$.
Let $Q$ be the inverse image of $M$ in $B$.
It is prime as an inverse image of a prime ideal.
For $r\in R$, its image lies in $M$ exactly when its image in $R_P$
lies in $PR_P$, and Theorem~\ref{thm:localization-primes} says the latter
condition is $r\in P$. Hence $Q\cap R=P$ as required.
\end{proof}
An injective extension is part of this statement. For a noninjective
integral algebra map the corresponding formulation applies to primes
containing the kernel, after factoring through its image.
The zero ring has no prime ideals, so causes no exception.
This bounded continuation establishes tensor and finiteness tools,
localization, and basic integrality with lying over. Full Groebner-basis
algorithms, the Nullstellensatz, Noether normalization, and ideal class
groups remain explicit topics for the advanced undergraduate D22
direction; they have not been proved here. Exterior and symmetric
powers over fields retain their precise D05 interface
\texttt{K05.15.exterior} and \texttt{K05.15.symmetric}.
Their general-ring and projective-module extensions are also deferred.

